\section{Lecture 26: Generative Adversarial Networks}

\subsection{Bridge: Explicit and Implicit Generative Models}

In the previous lecture, we studied a likelihood-based generative model:
\[
p_\theta(x)=\int p_\theta(x|z)p(z)\,dz.
\]

There, the model tries to assign probability to data explicitly.

\medskip

\noindent
\textbf{Explicit generative models.}
These define a density or likelihood, at least in principle:
\[
p_\theta(x).
\]

Examples include:
\begin{itemize}
    \item autoregressive models
    \item latent-variable models such as VAEs
    \item normalizing flows
\end{itemize}

\medskip

\noindent
There is another route.

\medskip

\noindent
\textbf{Implicit generative models.}
Instead of explicitly writing down a tractable density $p_\theta(x)$, we define a sampling procedure:
\[
z \sim p(z), \qquad x = G_\theta(z).
\]

Here:
\begin{itemize}
    \item $z$ is latent noise
    \item $G_\theta$ is a generator network
    \item generated samples induce a model distribution $p_g$
\end{itemize}

\medskip

\noindent
\textbf{Key difference.}
An implicit model may be easy to sample from but hard to evaluate as a density.

\medskip

\noindent
This raises the question:
\begin{quote}
If we cannot easily compute likelihood, how can we train a generative model?
\end{quote}

\medskip

\noindent
Generative adversarial networks answer this by turning learning into a game between two networks.

\subsection{Core Idea of GANs}

A generative adversarial network (GAN) consists of two parts:

\begin{itemize}
    \item a \textbf{generator} $G_\theta$
    \item a \textbf{discriminator} $D_\phi$
\end{itemize}

\medskip

\noindent
\textbf{Generator.}
The generator maps random noise to synthetic samples:
\[
z \sim p(z), \qquad \tilde{x}=G_\theta(z).
\]

\medskip

\noindent
\textbf{Discriminator.}
The discriminator receives an input $x$ and outputs a number in $(0,1)$:
\[
D_\phi(x)\in(0,1),
\]
interpreted as the probability that $x$ is real rather than generated.

\medskip

\noindent
\textbf{Adversarial setup.}

\begin{itemize}
    \item the discriminator tries to distinguish real data from fake samples
    \item the generator tries to produce samples that fool the discriminator
\end{itemize}

\medskip

\noindent
\textbf{Intuition.}
The discriminator acts as a learned critic. It does not merely classify data; it provides a moving training signal that tells the generator how unrealistic its outputs still are.

\subsection{The Minimax Objective}

Let $p_{\text{data}}$ denote the true data distribution and let $p_g$ denote the distribution induced by the generator.

\medskip

\noindent
The original GAN objective is
\[
\min_G \max_D V(D,G),
\]
where
\[
V(D,G)
=
\mathbb{E}_{x\sim p_{\text{data}}}[\log D(x)]
+
\mathbb{E}_{z\sim p(z)}[\log(1-D(G(z)))].
\]

\medskip

\noindent
\textbf{Meaning of each term.}

\begin{itemize}
    \item $\mathbb{E}_{x\sim p_{\text{data}}}[\log D(x)]$ rewards the discriminator for assigning high probability to real data
    \item $\mathbb{E}_{z\sim p(z)}[\log(1-D(G(z)))]$ rewards the discriminator for assigning low probability to fake data
\end{itemize}

\medskip

\noindent
The generator wants the opposite outcome: it wants generated samples to look real to the discriminator.

\medskip

\noindent
\textbf{Game interpretation.}
The discriminator tries to maximize classification success; the generator tries to minimize it by producing more realistic samples.

\subsection{Two-Player Optimization View}

Training a GAN is not ordinary minimization of one scalar loss over one set of parameters. It is a two-player game.

\medskip

\noindent
\textbf{Discriminator step.}
For fixed generator $G$, update $\phi$ to maximize
\[
V(D,G).
\]

\medskip

\noindent
\textbf{Generator step.}
For fixed discriminator $D$, update $\theta$ to minimize
\[
V(D,G).
\]

\medskip

\noindent
In practice, one alternates gradient-based updates:
\begin{enumerate}
    \item sample real data $x \sim p_{\text{data}}$
    \item sample noise $z \sim p(z)$ and form fake data $\tilde{x}=G(z)$
    \item update discriminator to improve real/fake separation
    \item update generator to make fake samples more convincing
\end{enumerate}

\medskip

\noindent
\textbf{Important difference from supervised learning.}
The target for the generator is not fixed in advance. It is created dynamically by the current state of the discriminator.

\subsection{Generator Loss Variants}

If the generator minimizes the original saturating objective
\[
\mathbb{E}_{z\sim p(z)}[\log(1-D(G(z)))],
\]
training can be slow when the discriminator is strong.

\medskip

\noindent
\textbf{Why?}
If $D(G(z))$ is already close to $0$, then $\log(1-D(G(z)))$ changes only weakly and the generator may receive poor gradients.

\medskip

\noindent
A common practical alternative is the \textbf{non-saturating} generator loss:
\[
\min_G -\mathbb{E}_{z\sim p(z)}[\log D(G(z))].
\]

\medskip

\noindent
\textbf{Interpretation.}
Instead of directly minimizing the discriminator's success, the generator is trained to maximize the discriminator's belief that fake samples are real.

\medskip

\noindent
This usually gives stronger gradients early in training.

\subsection{Optimal Discriminator Intuition}

To understand GANs theoretically, fix the generator $G$, hence fix the model distribution $p_g$.

\medskip

\noindent
For a given input $x$, the discriminator seeks to maximize
\[
p_{\text{data}}(x)\log D(x) + p_g(x)\log(1-D(x)).
\]

Pointwise optimization yields the optimal discriminator:
\[
D^*(x)=\frac{p_{\text{data}}(x)}{p_{\text{data}}(x)+p_g(x)}.
\]

\medskip

\noindent
\textbf{Interpretation.}
At each point $x$:
\begin{itemize}
    \item if real data density is much larger than fake density, then $D^*(x)$ is near $1$
    \item if fake density is much larger than real density, then $D^*(x)$ is near $0$
    \item if they match, then $D^*(x)=1/2$
\end{itemize}

\medskip

\noindent
\textbf{Key lesson.}
The optimal discriminator measures how distinguishable the two distributions are.

\subsection{What Happens at Equilibrium?}

Substituting the optimal discriminator back into the value function gives
\[
V(D^*,G)
=
\mathbb{E}_{x\sim p_{\text{data}}}
\left[
\log \frac{p_{\text{data}}(x)}{p_{\text{data}}(x)+p_g(x)}
\right]
+
\mathbb{E}_{x\sim p_g}
\left[
\log \frac{p_g(x)}{p_{\text{data}}(x)+p_g(x)}
\right].
\]

This can be simplified into
\[
V(D^*,G)
=
- \log 4 + 2\,\mathrm{JS}(p_{\text{data}}\|p_g),
\]
where $\mathrm{JS}$ denotes the Jensen--Shannon divergence.

\medskip

\noindent
\textbf{Consequence.}
If the discriminator is always optimal, GAN training pushes the generator toward minimizing a divergence between the real and generated distributions.

\medskip

\noindent
\textbf{Global optimum.}
The minimum occurs when
\[
p_g = p_{\text{data}},
\]
in which case
\[
D^*(x)=\frac12
\]
everywhere, since real and fake samples are no longer distinguishable.

\medskip

\noindent
\textbf{Insight.}
The discriminator is not the final goal. It is a device for guiding the generator toward the data distribution.

\subsection{Divergence Viewpoint}

GANs can be understood as learning by minimizing a discrepancy between distributions.

\medskip

\noindent
\textbf{In the original theory,} this discrepancy is the Jensen--Shannon divergence:
\[
\mathrm{JS}(p_{\text{data}}\|p_g).
\]

\medskip

\noindent
\textbf{Why this matters.}
It connects adversarial learning to a broader theme in statistics and machine learning:
\begin{quote}
Learn a model distribution by making it close to the true distribution under some notion of distance or divergence.
\end{quote}

\medskip

\noindent
Different generative methods use different notions of discrepancy:
\begin{itemize}
    \item maximum likelihood often corresponds to KL-related objectives
    \item VAEs use an ELBO involving likelihood and KL terms
    \item original GANs implicitly target Jensen--Shannon divergence
    \item Wasserstein GANs use Earth-Mover distance intuition
\end{itemize}

\medskip

\noindent
\textbf{Important practical caveat.}
A divergence may behave poorly when the supports of two distributions barely overlap, which is one reason GAN training can become unstable.

\subsection{Geometric Intuition}

Suppose the data lie on a thin low-dimensional manifold inside a high-dimensional ambient space. Early in training, the generator may place mass on a very different manifold.

\medskip

\noindent
Then:
\begin{itemize}
    \item the discriminator can separate real and fake samples almost perfectly
    \item gradients for the generator may become weak or uninformative
    \item optimization becomes unstable
\end{itemize}

\medskip

\noindent
\textbf{Intuition.}
In high dimensions, distributions can be visually similar in sample quality yet still be easy to separate statistically. This makes adversarial training delicate.

\subsection{Why GANs Can Produce Sharp Samples}

Compared with likelihood-based models such as standard VAEs, GANs are often associated with visually sharp outputs.

\medskip

\noindent
\textbf{Reason.}
The discriminator does not reward average-case pixelwise accuracy. Instead, it rewards realism under a learned criterion.

\medskip

\noindent
This can reduce the tendency to average over multiple plausible outputs.

\medskip

\noindent
\textbf{Tradeoff.}
GANs may produce highly realistic samples without giving a tractable density model, but the optimization is typically more fragile.

\subsection{Mode Collapse}

One of the most famous GAN pathologies is \textbf{mode collapse}.

\medskip

\noindent
\textbf{Definition.}
Mode collapse occurs when the generator maps many latent inputs $z$ to a small set of outputs or to only a few modes of the true data distribution.

\medskip

\noindent
Instead of covering all major patterns in the data, the generator may produce:
\begin{itemize}
    \item only a narrow family of faces
    \item only certain object poses
    \item repeated high-quality but low-diversity samples
\end{itemize}

\medskip

\noindent
\textbf{Why it can happen.}
If a small subset of outputs consistently fools the current discriminator, the generator may keep exploiting those outputs rather than learning the full distribution.

\medskip

\noindent
\textbf{Game-theoretic view.}
The generator may find a locally successful strategy against the present discriminator, even though that strategy is globally poor as a generative model.

\subsection{Training Instability}

GAN training is well known to be unstable.

\medskip

\noindent
\textbf{Common symptoms.}
\begin{itemize}
    \item discriminator quickly becomes too strong
    \item generator gradients vanish or oscillate
    \item sample quality improves and then degrades
    \item different runs behave very differently
\end{itemize}

\medskip

\noindent
\textbf{Why this is hard.}
Ordinary gradient descent is designed for minimization, but GANs involve a minimax game:
\[
\min_\theta \max_\phi V(\phi,\theta).
\]

Even simple games can produce cycling rather than convergence.

\medskip

\noindent
\textbf{Insight.}
GAN training difficulty is not just a matter of architecture or hyperparameters. It is tied to the geometry of adversarial optimization itself.

\subsection{Stabilization Tricks in Practice}

A large part of GAN research has focused on making training more stable.

\medskip

\noindent
\textbf{Common tricks include:}

\begin{itemize}
    \item \textbf{non-saturating generator loss} rather than the original saturating form
    \item \textbf{careful balance of updates} between generator and discriminator
    \item \textbf{label smoothing} for real targets
    \item \textbf{noise injection} into inputs or intermediate features
    \item \textbf{batch normalization} or architectural normalization methods
    \item \textbf{spectral normalization} to control discriminator Lipschitz behavior
    \item \textbf{feature matching} losses that encourage broader distributional alignment
    \item \textbf{minibatch discrimination} to reduce mode collapse
\end{itemize}

\medskip

\noindent
\textbf{Architectural design also matters.}
Convolutional inductive bias, normalization choices, and discriminator capacity strongly affect training behavior.

\medskip

\noindent
\textbf{Important principle.}
The discriminator should be strong enough to provide useful feedback, but not so overpowering that the generator receives no meaningful gradient signal.

\subsection{Conditional GANs}

Often we do not want unconditional generation. Instead, we want to generate data conditioned on side information.

\medskip

\noindent
Let $y$ denote a conditioning variable, such as:
\begin{itemize}
    \item a class label
    \item text
    \item another image
    \item semantic attributes
\end{itemize}

\medskip

\noindent
A conditional GAN models
\[
p(x|y).
\]

\medskip

\noindent
\textbf{Generator.}
\[
\tilde{x}=G(z,y).
\]

\medskip

\noindent
\textbf{Discriminator.}
\[
D(x,y),
\]
which judges whether $x$ is a real sample consistent with condition $y$.

\medskip

\noindent
\textbf{Conditional objective.}
\[
\min_G \max_D
\;
\mathbb{E}_{(x,y)\sim p_{\text{data}}}[\log D(x,y)]
+
\mathbb{E}_{z\sim p(z),\, y\sim p(y)}[\log(1-D(G(z,y),y))].
\]

\medskip

\noindent
\textbf{Applications.}
\begin{itemize}
    \item class-conditional image generation
    \item image-to-image translation
    \item text-to-image generation in early adversarial systems
    \item super-resolution
\end{itemize}

\medskip

\noindent
\textbf{Conceptual point.}
Conditioning reduces ambiguity and gives the generator a clearer target. Instead of modeling the entire distribution at once, it models a family of conditional distributions.

\subsection{Image-to-Image and Structured Conditional Generation}

Conditional GANs became especially influential in structured generation tasks.

\medskip

\noindent
Examples include:
\begin{itemize}
    \item edge map $\rightarrow$ photo
    \item segmentation mask $\rightarrow$ street scene
    \item low-resolution image $\rightarrow$ high-resolution image
\end{itemize}

\medskip

\noindent
In such cases, the generator is not producing from pure noise alone. It is transforming one representation into another while the discriminator enforces realism.

\medskip

\noindent
\textbf{Practical advantage.}
The adversarial loss can encourage sharper, more realistic outputs than purely pixelwise losses.

\subsection{Wasserstein GAN: Motivation}

Original GANs can suffer when Jensen--Shannon divergence becomes uninformative, especially if the real and generated distributions are far apart or supported on different low-dimensional sets.

\medskip

\noindent
Wasserstein GAN (WGAN) changes the discrepancy measure.

\medskip

\noindent
\textbf{Idea.}
Instead of using a discriminator that estimates a classification probability, use a critic that scores samples, and train the generator to reduce the \emph{Wasserstein-1 distance} between distributions.

\medskip

\noindent
This distance is also called the \textbf{Earth-Mover distance}.

\medskip

\noindent
\textbf{Intuition.}
Imagine one distribution as a pile of earth and the other as a desired arrangement. The Wasserstein distance measures how much ``work'' is needed to move mass from one to the other.

\medskip

\noindent
\textbf{Why this helps.}
Unlike some divergences, Wasserstein distance can still provide meaningful gradients even when distributions have little overlap.

\subsection{Conceptual WGAN Objective}

In WGAN, the discriminator is replaced by a \textbf{critic} $f_\phi(x)$, which outputs a real number rather than a probability.

\medskip

\noindent
Conceptually, the objective takes the form
\[
\max_{f \in \mathcal{F}_{\text{Lip}}}
\;
\mathbb{E}_{x\sim p_{\text{data}}}[f(x)]
-
\mathbb{E}_{x\sim p_g}[f(x)],
\]
where $\mathcal{F}_{\text{Lip}}$ denotes a class of 1-Lipschitz functions.

\medskip

\noindent
The generator then tries to minimize this gap.

\medskip

\noindent
\textbf{Interpretation.}
The critic assigns higher scores to real samples and lower scores to fake ones, under a smoothness constraint. This score difference estimates how far the two distributions are in Wasserstein distance.

\subsection{Why the Lipschitz Constraint Matters}

The Lipschitz constraint is central to WGAN theory.

\medskip

\noindent
\textbf{Without it,} the critic could scale arbitrarily and the Wasserstein interpretation would break.

\medskip

\noindent
\textbf{With it,} the critic remains controlled and yields a meaningful notion of distance between distributions.

\medskip

\noindent
The original WGAN enforced this approximately by \textbf{weight clipping}, but this can damage optimization and restrict critic capacity.

\subsection{WGAN-GP: Gradient Penalty}

A more effective alternative is \textbf{WGAN-GP}, where the Lipschitz constraint is encouraged by penalizing gradient norm deviations.

\medskip

\noindent
Conceptually, one adds a term of the form
\[
\lambda \, \mathbb{E}_{\hat{x}}
\left[
\bigl(\|\nabla_{\hat{x}} f(\hat{x})\|_2 - 1\bigr)^2
\right],
\]
where $\hat{x}$ is sampled along lines between real and fake examples.

\medskip

\noindent
\textbf{Interpretation.}
This encourages the critic to behave approximately like a 1-Lipschitz function.

\medskip

\noindent
\textbf{Practical outcome.}
WGAN-GP often leads to:
\begin{itemize}
    \item more stable training
    \item more meaningful loss curves
    \item improved resistance to severe gradient pathologies
\end{itemize}

\subsection{Comparing Original GAN and WGAN}

\textbf{Original GAN.}
\begin{itemize}
    \item discriminator outputs a probability
    \item linked to Jensen--Shannon divergence
    \item can suffer from saturation and instability
\end{itemize}

\medskip

\noindent
\textbf{WGAN.}
\begin{itemize}
    \item critic outputs a real-valued score
    \item linked to Wasserstein distance intuition
    \item often provides smoother optimization behavior
\end{itemize}

\medskip

\noindent
\textbf{Conceptual lesson.}
Changing the notion of discrepancy between distributions changes the geometry of the learning problem.

\subsection{GANs in the Broader Generative Modeling Landscape}

GANs occupy a distinct point in the design space of generative models.

\medskip

\noindent
\textbf{Compared with VAEs:}
\begin{itemize}
    \item GANs often produce sharper samples
    \item VAEs offer a more explicit probabilistic framework
    \item GANs avoid explicit likelihood
    \item VAEs provide a clearer inference model through the encoder
\end{itemize}

\medskip

\noindent
\textbf{Compared with diffusion models:}
\begin{itemize}
    \item GAN generation is often fast at sampling time
    \item diffusion models are usually more stable to train
    \item diffusion models provide broader mode coverage in many modern settings
\end{itemize}

\medskip

\noindent
\textbf{Historical role.}
GANs were crucial in showing that high-dimensional neural generators could produce strikingly realistic samples, and they shaped much of modern generative modeling.

\subsection{Summary}

\textbf{Main points of this lecture.}

\begin{itemize}
    \item GANs are implicit generative models defined by a sampling procedure
    \[
    z\sim p(z), \qquad x=G(z).
    \]
    \item Learning is framed as an adversarial game between generator and discriminator.
    \item The original minimax objective is
    \[
    \min_G \max_D
    \left[
    \mathbb{E}_{x\sim p_{\text{data}}}[\log D(x)]
    +
    \mathbb{E}_{z\sim p(z)}[\log(1-D(G(z)))]
    \right].
    \]
    \item For fixed generator, the optimal discriminator is
    \[
    D^*(x)=\frac{p_{\text{data}}(x)}{p_{\text{data}}(x)+p_g(x)}.
    \]
    \item Under this discriminator, GAN training relates to minimizing Jensen--Shannon divergence.
    \item GANs can produce sharp samples but are prone to instability and mode collapse.
    \item Stabilization methods include improved losses, normalization, regularization, and architectural choices.
    \item Conditional GANs generate data conditioned on labels or other inputs.
    \item WGAN and WGAN-GP replace the original divergence picture with Wasserstein distance intuition and often improve training behavior.
\end{itemize}

\medskip

\noindent
\textbf{Looking ahead.}
GANs generate in one shot through adversarial training. In the next lecture, we study diffusion models, which generate instead by gradual denoising and iterative refinement.